\documentclass[10pt,a4paper]{article}
\usepackage{lmodern}
\usepackage[margin=18mm]{geometry}
\usepackage{booktabs}
\usepackage{graphicx}
\usepackage[table]{xcolor}
\usepackage[hidelinks]{hyperref}
\usepackage{parskip}
\usepackage{array}
\usepackage{enumitem}
\graphicspath{{figures/}}
\definecolor{warning}{HTML}{FFF3CD}
\definecolor{info}{HTML}{E0F2FE}
\definecolor{good}{HTML}{DCFCE7}
\definecolor{bad}{HTML}{FEE2E2}
\hypersetup{
  pdftitle={Dynamic and static UWB/IMU validation -- 11 August 2026},
  pdfauthor={UWB ESP-IDF technical report}
}
\setlist{nosep}
\begin{document}

\begin{titlepage}
\centering
\vspace*{16mm}
{\Huge\bfseries Dynamic and static UWB/IMU validation\par}
\vspace{4mm}
{\Large FlexTDOA, Native DS-TWR and Passive DS-TWR\par}
\vspace{11mm}
{\large Field campaign and controlled replay ablation -- 11 August 2026\par}
\vspace{12mm}
\begin{tabular}{rl}
Platform: & 5 ESP32-S3 + DW3000 modules; Raspberry Pi collector\\
Topology: & M1 tag; M2, M3, M4 and M5 anchors\\
Radio: & UWB channel 9; configured 40 MHz SPI\\
Firmware: & deployed \texttt{bb23b07-dirty}; analysis branch \texttt{imu-fusion-smoothing}\\
Sensor path: & BNO085 acquisition near 500 Hz; GPTimer timestamps\\
Evidence: & 3 dynamic + 3 static protocol captures; replay and ablation artifacts\\
\end{tabular}
\vfill
\colorbox{warning}{\parbox{0.89\textwidth}{\vspace{2mm}
\textbf{Main result.} Raw UWB remains the authoritative positioning stream.
The current fusion output can suppress isolated extremes, but controlled
replays show that the measurable improvement comes from the common position
smoother, not from accelerometer integration. The accelerometer contribution
is therefore \textbf{not accepted} for production accuracy claims.
\vspace{2mm}}}
\vfill
{\large Reproducible technical report\par}
\end{titlepage}

\pagenumbering{roman}
\tableofcontents
\clearpage
\pagenumbering{arabic}

\section{Executive summary}

This report is the complete counterpart of the 6 August protocol report. It
uses the same six-cell matrix (three protocols in dynamic and static modes),
the same independent-frame policy, the same local static precision metrics,
and the same RTK registration audit. It additionally contains the controlled
IMU ablation that is necessary to distinguish inertial information from
ordinary trajectory smoothing.

\colorbox{info}{\parbox{0.96\textwidth}{\textbf{Evidence boundary.}
Counts, rates, positive gaps and local static clouds are direct diagnostics.
Dynamic UWB--RTK errors are useful in-system comparisons, but the UWB packet
does not yet carry the radio measurement timestamp. Static absolute ranking is
invalid for the original three 30 s captures because their RTK centers differ
by 2.03 m and two anchor registrations are poor. No outlier was removed.}}

The principal results are:
\begin{itemize}
\item Static raw local RMS is \textbf{2.13 cm} for FlexTDOA,
      \textbf{2.02 cm} for Native DS-TWR and \textbf{1.78 cm} for Passive
      DS-TWR. These values measure repeatability around each capture median.
\item Dynamic independent rates are \textbf{12.90}, \textbf{20.86} and
      \textbf{13.33 positions/s}, respectively. Flex and Passive contain
      multi-second maximum gaps; Native contains a 1.02 s maximum gap.
\item On the controlled Native replay, raw RMSE is 17.911 cm. Disabling
      acceleration gives 16.323 cm; the current acceleration path gives
      16.326 cm. The difference is \textbf{+0.0031 cm}: acceleration provides
      no measurable improvement.
\item Passive gives an even stronger null result: disabled, current and
      relaxed acceleration configurations all produce 14.867 cm RMSE and
      29.389 cm P95 to the displayed precision.
\item A new Passive stationary control with the full 500 Hz telemetry path
      yields \textbf{2.710 cm raw RMSE versus RTK} and \textbf{1.711 cm local
      raw RMS}. The radio/solver path has therefore not suffered a general
      stationary regression from enabling high-rate IMU acquisition.
\end{itemize}

\section{Questions and acceptance criteria}

The campaign evaluates four separate questions:
\begin{enumerate}
\item Does each protocol continue to deliver raw positions at a useful rate?
\item What is the unfiltered stationary repeatability of each protocol?
\item Does the current fused output improve dynamic RMSE and P95 without
      increasing overshoot or hiding raw data?
\item Is any observed improvement caused by acceleration, rather than by the
      position smoother shared by all replay variants?
\end{enumerate}

The causal acceptance criterion is intentionally strict: acceleration must
beat an otherwise identical replay with acceleration disabled on both RMSE and
P95, without a worse maximum or new dead zones. A raw-versus-fused comparison
alone cannot prove an IMU gain because the fused branch also includes endpoint
smoothing, constant-velocity propagation, reacquisition logic and ZUPT.

\section{System and protocol background}

\subsection{FlexTDOA}
The tag is UWB receive-only. Anchors exchange scheduled transmissions and the
tag solves from a coherent set of time-difference observations. Airtime scales
well with passive tag count. The current capture includes both independent
frames and repeated/publication events; only independent frames enter fair
rate and trajectory metrics.

\subsection{Native DS-TWR}
The active tag ranges to every anchor and solves one coherent four-range frame.
It offers the clearest measurement semantics but consumes tag airtime. A frame
can be incomplete when one or more anchor exchanges miss their deadline.

\subsection{Passive DS-TWR}
The tag listens to anchor-to-anchor DS-TWR exchanges. Rolling windows can
publish multiple positions from shared measurements. Therefore all-event rate
and independent-frame rate are reported separately.

\subsection{IMU path under test}
BNO085 samples carry 64-bit GPTimer time into replay. Orientation association,
reordering and interpolation are available, and raw UWB remains permanently
separate. The filter can use constant-velocity propagation, position endpoint
smoothing, stationary bias estimation, ZUPT and yaw-conditioned horizontal
acceleration. That mixture is the reason for the controlled ablation.

\section{Captured evidence}

The selected dynamic walks lasted 177--243 s; the static blocks lasted about
30 s. The Passive dynamic JSONL lacks its final \texttt{capture\_end} record,
but its 177 s data span and replay are readable. It is included and explicitly
marked incomplete rather than silently discarded.

\begin{center}\scriptsize\resizebox{\textwidth}{!}{%
\begin{tabular}{lllrrrrrl}\toprule
\textbf{Protocol} & \textbf{Mode} & \textbf{capture ID} & \textbf{duration s} &
\textbf{events} & \textbf{independent} & \textbf{indep Hz} & \textbf{coverage \%} & \textbf{container}\\
\midrule
FlexTDOA & dynamic & flextdoa\_dynamic\_final\_01 & 240.5 & 4023 & 3101 & 12.90 & 99.4 & complete\\
Native DS-TWR & dynamic & native\_ds\_dynamic\_quality\_gate\_02 & 243.0 & 5069 & 5069 & 20.86 & 99.7 & complete\\
Passive DS-TWR & dynamic & passive\_ds\_dynamic\_final\_01 & 176.9 & 7033 & 2358 & 13.33 & 100.0 & missing capture\_end\\
FlexTDOA & static & flextdoa\_static\_final & 30.3 & 929 & 929 & 30.67 & 99.4 & complete\\
Native DS-TWR & static & native\_ds\_static\_final & 30.5 & 686 & 686 & 22.47 & 99.4 & complete\\
Passive DS-TWR & static & passive\_ds\_static\_final & 31.1 & 2295 & 772 & 24.80 & 99.1 & complete\\
\bottomrule\end{tabular}}
\end{center}

\section{Methodology}

\subsection{Population and continuity}
Interleaved records marked as fused are excluded from raw UWB metrics. Passive
and Flex publications are reduced to independent frames where the protocol
marks them. Rates use the full capture wall duration. Positive gaps are
differences between unique ESP32 uptime values; P50, P95, P99, maximum and
counts above 100 ms are retained.

\subsection{Static precision}
Raw radial displacement is calculated about the component-wise median of each
capture. CEP50, RMS, P95, P99, maximum, 2DRMS and first-to-last-decile drift
retain every independent sample. These metrics describe local repeatability,
not absolute accuracy.

\subsection{RTK association and frame registration}
Only quality-4 RTK Fixed fixes enter comparisons. WGS84 is converted to ENU. A
proper 2-D rigid rotation and translation (no scale) is fitted from anchor
geometry, never from the tag trajectory. UWB and RTK samples are associated in
software within 100 ms.

The measurement-time estimate currently subtracts reported GNSS fix age from
collector receive time. A separate audit found this estimate about 108--112 ms
late relative to GNSS UTC; correcting it explains roughly 5--6 cm of dynamic
RMSE for Flex and Native, but not the complete 15--18 cm error. The radio
measurement timestamp must be transmitted end-to-end before sub-decimetre
dynamic claims are treated as final.

\subsection{Controlled IMU ablation}
Raw, acceleration-disabled, current-gate and relaxed-gate configurations are
replayed on the same captured population. The disabled replay preserves the
same position smoother and constant-velocity machinery. Consequently,
\emph{current minus disabled} isolates the acceleration contribution much more
cleanly than \emph{fused minus raw}.

\section{Dynamic raw results}

\subsection{Trajectories}
The walks are separate physical trials rather than synchronized repetitions.
The panels are therefore used for continuity, jumps and gross geometry, not
point-by-point route ranking.

\begin{figure}[p]
\centering
\includegraphics[width=\textwidth]{06_dynamic_trajectories.pdf}
\caption{Raw independent dynamic trajectories in each local UWB frame. Green
markers are anchors. Outliers are retained.}
\end{figure}

\subsection{Update rate and gaps}

\begin{center}\scriptsize\resizebox{\textwidth}{!}{%
\begin{tabular}{lrrrrrrrr}\toprule
\textbf{Protocol} & \textbf{event Hz} & \textbf{indep Hz} & \textbf{P50 ms} &
\textbf{P95 ms} & \textbf{P99 ms} & \textbf{max ms} & \textbf{$>$100 ms} & \textbf{step P95 m/s}\\
\midrule
FlexTDOA & 16.73 & 12.90 & 30 & 160 & 280 & 2550 & 380 & 14.75\\
Native DS-TWR & 20.86 & 20.86 & 40 & 70 & 130 & 1020 & 71 & 5.22\\
Passive DS-TWR & 39.75 & 13.33 & 10 & 120 & 240 & 2680 & 385 & 4.29\\
\bottomrule\end{tabular}}
\end{center}

\begin{figure}[htbp]
\centering
\includegraphics[width=0.94\textwidth]{01_position_rate.pdf}
\caption{Independent position rate. Passive all-event publication rate is not
substituted for independent frames.}
\end{figure}

\begin{figure}[htbp]
\centering
\includegraphics[width=0.94\textwidth]{02_position_gap.pdf}
\caption{Positive-gap P95 for dynamic and static raw streams. Maximum gaps are
listed in the tables.}
\end{figure}

Native has the strongest dynamic cadence in this campaign. Flex and Passive
both contain multi-second interruptions. The very large Flex step-speed P95 is
a direct symptom of abrupt raw position variation, not measured walking speed.

\subsection{Native radio-loss diagnostic}
Across the Native capture, start/end counter deltas give 4254 complete and
1208 incomplete candidate frames: \textbf{22.12\% incomplete}. There were 658
slot overruns. All 485 PHY errors were reported as recovered, but the same
interval also contained 964 response timeouts, 517 result timeouts, 377 frame
sync losses, 90 PHR errors and 18 FCS errors.

\begin{center}\scriptsize
\begin{tabular}{lrrrrr}\toprule
\textbf{Anchor} & \textbf{attempts} & \textbf{completed ranges} &
\textbf{completion \%} & \textbf{response timeout} & \textbf{result timeout}\\
\midrule
A2 & 5462 & 5028 & 92.06 & 302 & 131\\
A3 & 5461 & 5042 & 92.33 & 289 & 130\\
A4 & 5461 & 5219 & 95.57 & 120 & 122\\
A5 & 5461 & 5074 & 92.91 & 253 & 134\\
\bottomrule\end{tabular}
\end{center}

Loss is not confined to one anchor, although A4 is visibly better. The data do
not support blaming geometry alone. Scheduling and RX recovery remain credible
contributors.

\section{Static raw results}

\subsection{Local precision}

\begin{center}\scriptsize\resizebox{\textwidth}{!}{%
\begin{tabular}{lrrrrrrrr}\toprule
\textbf{Protocol} & \textbf{N} & \textbf{CEP50 cm} & \textbf{RMS cm} &
\textbf{P95 cm} & \textbf{P99 cm} & \textbf{max cm} & \textbf{2DRMS cm} & \textbf{drift cm}\\
\midrule
FlexTDOA & 929 & 1.68 & 2.13 & 3.85 & 4.89 & 5.95 & 4.27 & 0.21\\
Native DS-TWR & 686 & 1.40 & 2.02 & 3.76 & 5.99 & 10.06 & 4.02 & 0.16\\
Passive DS-TWR & 772 & 1.50 & 1.78 & 3.11 & 3.72 & 6.44 & 3.57 & 0.45\\
\bottomrule\end{tabular}}
\end{center}

\begin{figure}[htbp]
\centering
\includegraphics[width=0.94\textwidth]{03_static_precision.pdf}
\caption{Static raw precision around each capture median. All independent
samples are included.}
\end{figure}

\begin{figure}[p]
\centering
\includegraphics[width=\textwidth]{07_static_position_clouds.pdf}
\caption{Static raw position clouds centered on each capture median and drawn
at one common scale. The cross is the median.}
\end{figure}

\begin{figure}[p]
\centering
\includegraphics[width=\textwidth]{08_static_gps_rtk_clouds.pdf}
\caption{Static GPS RTK Fixed tag solutions. Top: each cloud centered on its
own median at a shared centimetre scale. Bottom: the same fixes in a common ENU
frame, exposing the cross-capture reference displacement.}
\end{figure}

All three protocols remain close to the requested 2 cm local RMS target at the
tested central point. This does not establish 2 cm absolute accuracy across the
polygon or under NLOS.

\subsection{Static RTK consistency gate}
The maximum separation between the three stationary RTK centers is 2.032 m,
versus a 0.250 m consistency gate. Flex has no usable anchor transform in its
static block. Native and Passive anchor-fit RMS values are 1.073 m and 0.741 m,
respectively. Therefore the approximately 3.0 m Native and 1.9 m Passive
absolute offsets produced by these blocks are registration artifacts and
\textbf{must not be used to rank protocols}.

\subsection{Current stationary regression control}
A new 30 s Passive capture made after the high-rate IMU changes has valid RTK
registration and provides the relevant regression check.

\begin{center}\small
\begin{tabular}{lr}\toprule
\textbf{Metric} & \textbf{result}\\
\midrule
Independent raw positions & 832\\
Raw UWB RMSE versus RTK & 2.710 cm\\
Raw UWB P95 versus RTK & 4.080 cm\\
Raw UWB maximum versus RTK & 4.993 cm\\
Raw local-cloud RMS & 1.711 cm\\
Raw local-cloud P95 & 2.915 cm\\
Fused RMSE versus RTK & 2.568 cm\\
Anchor rigid-fit RMS & 0.422 cm\\
RTK rate & 8.08 Hz\\
Tag IMU rate used by replay & 458.9 Hz\\
\bottomrule\end{tabular}
\end{center}

This is close to the prior calibrated stationary baseline of 2.55 cm absolute
RMSE. High-rate sensor acquisition and telemetry have not generally degraded
stationary raw UWB accuracy.

\section{Dynamic UWB--RTK audit}

\begin{center}\scriptsize\resizebox{\textwidth}{!}{%
\begin{tabular}{lrrrrrrrr}\toprule
\textbf{Protocol} & \textbf{pairs} & \textbf{pair coverage \%} &
\textbf{time P95 ms} & \textbf{raw RMSE cm} & \textbf{raw P95 cm} &
\textbf{raw max cm} & \textbf{anchor fit P95 cm} & \textbf{debiased RMSE cm}\\
\midrule
FlexTDOA & 1100 & 62.6 & 65.6 & 18.28 & 30.57 & 80.19 & 1.41 & 17.16\\
Native DS-TWR & 1892 & 98.2 & 28.0 & 18.68 & 28.36 & 278.00 & 2.25 & 14.77\\
Passive DS-TWR & 1154 & 83.3 & 83.4 & 16.75 & 34.66 & 67.91 & 1.50 & 16.82\\
\bottomrule\end{tabular}}
\end{center}

These values use direct capture-time association and anchor-based rigid
registration. Replay tables later use their own identical-timestamp population;
small differences in sample selection are expected and are retained rather
than forced to agree.

\begin{figure}[htbp]
\centering
\includegraphics[width=0.94\textwidth]{04_dynamic_rtk_error_cdf.pdf}
\caption{Dynamic raw UWB--RTK disagreement CDF after anchor registration. RTK
is an in-system reference, not surveyed truth.}
\end{figure}

\begin{figure}[htbp]
\centering
\includegraphics[width=0.94\textwidth]{05_alignment_quality.pdf}
\caption{Anchor-registration residual beside paired tag disagreement. Dynamic
anchor fits are centimetric; timing remains the larger limitation.}
\end{figure}

\section{Fusion replay and causal ablation}

\subsection{Raw versus selected fused replay}
The table below uses the per-tag comparison population and the selected replay
for each protocol. It describes the complete fused algorithm, not acceleration
alone.

\begin{center}\scriptsize\resizebox{\textwidth}{!}{%
\begin{tabular}{lrrrrrrrrl}\toprule
\textbf{Protocol} & \textbf{pairs} & \textbf{raw RMSE cm} & \textbf{fused RMSE cm} &
\textbf{raw P95 cm} & \textbf{fused P95 cm} & \textbf{raw max cm} &
\textbf{fused max cm} & \textbf{gap count delta} & \textbf{gate}\\
\midrule
FlexTDOA & 1044 & 17.262 & 16.883 & 28.062 & 28.263 & 78.42 & 65.12 & -190 & reject\\
Native DS-TWR & 1860 & 17.911 & 16.326 & 26.822 & 25.395 & 260.87 & 103.49 & -57 & nominal pass\\
Passive DS-TWR & 1067 & 14.953 & 15.029 & 29.440 & 29.818 & 67.78 & 67.41 & -355 & reject\\
\bottomrule\end{tabular}}
\end{center}

Flex reduces RMSE and maximum but increases P95 by 2.0 mm. Passive slightly
worsens both RMSE and P95. Native appears to pass a raw-versus-fused gate, but
that result is not causal evidence for the accelerometer.

\subsection{Controlled acceleration ablation}

\begin{center}\scriptsize\resizebox{\textwidth}{!}{%
\begin{tabular}{llrrrrr}\toprule
\textbf{Protocol} & \textbf{configuration} & \textbf{accel used / valid} &
\textbf{RMSE cm} & \textbf{P95 cm} & \textbf{max cm} & \textbf{generic replay pass}\\
\midrule
Native DS-TWR & raw UWB & -- & 17.910850 & 26.821895 & 260.873390 & --\\
Native DS-TWR & acceleration disabled & 4 / 108563 & 16.323327 & 25.395069 & 105.326516 & yes\\
Native DS-TWR & current gates & 20449 / 108563 & 16.326430 & 25.395069 & 103.493998 & yes\\
Native DS-TWR & relaxed gates & 97763 / 108563 & 16.310498 & 25.649909 & 94.720490 & yes\\
\midrule
Passive DS-TWR & raw UWB & -- & 14.952546 & 29.439575 & 67.783143 & --\\
Passive DS-TWR & acceleration disabled & 0 / 78893 & 14.867235 & 29.388807 & 66.904426 & yes\\
Passive DS-TWR & current gates & 0 / 78893 & 14.867235 & 29.388807 & 66.904426 & yes\\
Passive DS-TWR & relaxed gates & 26903 / 78893 & 14.867235 & 29.388807 & 66.904426 & yes\\
\bottomrule\end{tabular}}
\end{center}

\begin{figure}[htbp]
\centering
\includegraphics[width=0.94\textwidth]{09_imu_ablation_rmse.pdf}
\caption{Dynamic RMSE under controlled replay ablation. The large raw-to-
disabled change belongs to the common smoother; current acceleration is
indistinguishable from disabled.}
\end{figure}

\begin{figure}[htbp]
\centering
\includegraphics[width=0.94\textwidth]{10_imu_ablation_p95.pdf}
\caption{Dynamic P95 under the same ablation. Relaxing Native acceleration
gates slightly worsens P95.}
\end{figure}

For Native, current minus disabled is +0.003103 cm RMSE and exactly 0.000 cm
P95 to the shown precision. Relaxed acceleration changes RMSE by only
-0.012829 cm while worsening P95 by 0.254840 cm. Passive is numerically
identical even after 26903 samples are allowed to propagate. The current
vector-acceleration path therefore fails the causal acceptance criterion.

\subsection{Why acceleration was often unavailable}

\begin{center}\small
\begin{tabular}{lrl}\toprule
\textbf{Protocol} & \textbf{propagated / valid} & \textbf{dominant exclusive gates}\\
\midrule
Native DS-TWR & 20449 / 108563 = 18.84\% & yaw stale 69.30\%; yaw absent 9.28\%\\
FlexTDOA & 0 / 98805 & yaw absent 39.20\%; yaw stale 34.87\%; bias missing 25.93\%\\
Passive DS-TWR & 0 / 78893 & yaw stale 56.92\%; yaw absent 25.53\%; bias missing 17.55\%\\
\bottomrule\end{tabular}
\end{center}

\begin{figure}[htbp]
\centering
\includegraphics[width=0.90\textwidth]{11_accel_propagation_coverage.pdf}
\caption{Current-gate acceleration propagation coverage. Zero does not mean
that IMU transport failed; it means vector acceleration was blocked.}
\end{figure}

Flex has an additional replay epoch defect: an IMU timer backstep is translated
into an artificial $2^{32}$ ms span and triggers a reset. Its zero-propagation
figure cannot be attributed only to the normal gates. The current orientation
mapping also assumes body +X follows travel, which is false for arbitrary
handheld rotation and for the simulated throw.

\section{Root-cause findings beyond the filter}

\subsection{Missing radio measurement timestamp}
The current position wire format contains a 32-bit \texttt{uptime\_ms} assigned
when telemetry is submitted, not the radio measurement time. Dashboard and
replay therefore fall back to receive time. The required correction is a
versioned position stream carrying a 64-bit
\texttt{measurement\_monotonic\_us}, derived from the midpoint of the radio RX
host timestamps and preserved through the dashboard archive.

\subsection{Scheduling risks on the UWB core}
The UWB task runs on core 1 at priority 20 and installs the global GPIO ISR
service. BNO085 starts later on core 0, but its 500 Hz handler can consequently
execute on core 1 and perform a GPTimer read plus a critical section. Telemetry
also runs on core 1; its dequeue path scans up to 1024 PSRAM ring entries under
a critical section. Task priority cannot preempt code while interrupts are
masked. Both paths can increase missed radio deadlines and should be A/B tested
after restoring an O(1) FIFO or separate priority queues.

\subsection{What the IMU can safely do now}
Until a physical body-to-UWB mount calibration and a stable yaw source exist,
the safe uses are stationary detection, bias observation, ZUPT, motion-state
classification and adaptive process noise. Vector acceleration should remain
in a shadow predictor and must not replace raw UWB positions.

\section{Protocol assessment}

\begin{center}\small\resizebox{\textwidth}{!}{%
\begin{tabular}{p{0.18\textwidth}p{0.25\textwidth}p{0.25\textwidth}p{0.25\textwidth}}\toprule
\textbf{Criterion} & \textbf{FlexTDOA} & \textbf{Native DS-TWR} & \textbf{Passive DS-TWR}\\
\midrule
Tag radio role & receive-only; scalable & active exchange with each anchor & receive-only\\
Dynamic independent rate & 12.90 Hz & best: 20.86 Hz & 13.33 Hz after overlap removal\\
Dynamic worst gap & 2.55 s & best: 1.02 s & 2.68 s\\
Static local RMS & 2.13 cm & 2.02 cm & best: 1.78 cm\\
Current fused replay & rejected: P95 worsens 2 mm & generic pass, but not an IMU gain & rejected on per-tag RMSE/P95\\
Operational role & primary scalable mode after timestamp fix & coherent diagnostic baseline & useful passive mode; continuity still matters\\
\bottomrule\end{tabular}}
\end{center}

No protocol rewrite is justified by these data. The largest immediate gains are
expected from correct measurement timing, reduced UWB-core interference and
radio-loss diagnostics, not from more aggressive inertial integration.

\section{Limitations}
\begin{itemize}
\item The three dynamic routes are separate walks, not one synchronized path.
\item RTK is an in-system comparison reference and uses software association.
\item GNSS and UWB antenna lever arms are not surveyed.
\item The original static blocks fail the common-reference consistency gate;
      only local precision is comparable across them.
\item One stationary point does not cover GDOP, polygon edges, NLOS or tag
      orientation.
\item The Passive dynamic container lacks \texttt{capture\_end}; its data are
      included but its completeness status is not promoted.
\item The causal ablation is available for Native and Passive. Flex requires a
      replay after the timer-epoch defect is corrected.
\end{itemize}

\section{Recommended next validation sequence}
\begin{enumerate}
\item Add and validate the 64-bit radio measurement timestamp end-to-end.
\item Replace telemetry's O(1024) critical-section scan with O(1) queues and
      A/B test IMU off, acquisition-only and full streaming.
\item Keep yaw-free IMU assistance limited to stationary detection, ZUPT and
      adaptive process noise.
\item Calibrate the physical IMU-to-UWB rotation before using vector
      acceleration; do not infer body +X from arbitrary walking direction.
\item Repeat one marked dynamic route for all protocols and collect at most
      30 s static at the same point without restarting the RTK reference.
\item Accept fusion only when current beats acceleration-disabled on held-out
      RMSE and P95, with no larger maximum or dead-zone count.
\item After timing is trustworthy, investigate channel-9 antenna delay and
      per-link tag--anchor biases to approach the absolute $\pm$2 cm target.
\end{enumerate}

\section{Final conclusion}

The campaign is a useful negative result. High-rate IMU transport works and
stationary raw UWB remains centimetric, but the accelerometer has not improved
dynamic positioning. The visually calmer fused trace is primarily a smoother.
Calling the 17.91 to 16.33 cm Native change an accelerometer improvement would
be incorrect: the acceleration-disabled replay already reaches 16.323 cm.

Raw UWB remains the production and audit reference. Fusion stays optional and
experimental until measurement timing, orientation calibration and causal
held-out tests demonstrate a real improvement.

\section{Reproduction and audit artifacts}

The report directory contains:
\begin{itemize}
\item \texttt{analysis\_summary.json}, \texttt{capture\_metrics.csv},
      \texttt{replay\_metrics.csv}, \texttt{rtk\_pairs.csv} and alignment CSVs;
\item \texttt{ablation\_metrics.csv} and the script that generates the
      additional ablation figures;
\item eleven SVG/PDF figures, this TeX source and the compiled PDF;
\item source paths to the immutable JSONL captures. Raw archives were not
      duplicated into this directory.
\end{itemize}

The six matrix captures are in
\texttt{reports/imu\_fusion\_final\_validation\_20260811} and
\texttt{reports/imu\_fusion\_static\_20260811}. The current Passive regression
control is
\texttt{passive\_ds\_static\_raw\_regression\_01.passive\_ds.jsonl}.

\begin{verbatim}
python tools/uwb_dynamic_static_report.py ... --include-incomplete
python reports/imu_fusion_ablation_20260811/generate_additional_figures.py
pdflatex -interaction=nonstopmode report.tex
\end{verbatim}

No raw capture is modified during report generation.

\end{document}
